%% ===========================================================================
%%  main.tex — Quantifying with Confidence: Conformal Prediction for
%%           Reliable Musical Style Similarity Ranking
%% ===========================================================================

\documentclass[11pt,a4paper]{article}

% --- Packages ---
\usepackage[utf8]{inputenc}
\usepackage[T1]{fontenc}
\usepackage{amsmath,amssymb}
\usepackage{graphicx}
\usepackage{booktabs}
\usepackage{multirow}
\usepackage{hyperref}
\usepackage[margin=1in]{geometry}
\usepackage{enumitem}
\usepackage{caption}
\usepackage{subcaption}
\usepackage{float}
\usepackage{cite}

% --- Metadata ---
\title{Quantifying with Confidence:\\
       Conformal Prediction for Reliable Musical Style\\
       Similarity Ranking}
\author{}
\date{}

\begin{document}

\maketitle

\begin{abstract}
Quantifying musical style similarity is a central challenge in music information
retrieval (MIR) and a critical component of AI-generated music evaluation.
Current approaches based on self-supervised audio embeddings, such as
MERT-v1-330M, produce point-estimate similarity scores with no accompanying
measure of uncertainty. We introduce a framework that integrates conformal
prediction (CP) with self-supervised audio representations to construct
prediction sets with rigorous distribution-free coverage guarantees. Our method
extracts 1024-dimensional embeddings from MERT-v1-330M, constructs per-genre
prototype centroids from a reference set, and calibrates nonconformity scores
on a held-out calibration set using cosine distance. On the GTZAN dataset
(999 tracks, 10 genres), CP achieves empirical coverages of $0.993 \pm 0.010$,
$0.955 \pm 0.021$, and $0.910 \pm 0.028$ at significance levels
$\alpha = 0.01$, $0.05$, and $0.10$, respectively, tightly matching the nominal
targets. On the Free Music Archive (FMA) Small dataset (7{,}994 tracks, 8
genres)---a larger, noisier, and more diverse collection where MERT embeddings
exhibit markedly weaker class separation ($\Delta = 0.026$ vs.\ GTZAN's
$0.044$)---CP achieves coverages of $0.988$, $0.948$, and $0.898$, confirming
that the distribution-free guarantee is representation-agnostic. In contrast,
a Gaussian approximation baseline undercovers at
$\alpha = 0.01$ ($0.965$--$0.966$ vs.\ target $0.990$) on both datasets,
confirming that perceptual
similarity scores exhibit heavy-tailed distributions. Bootstrap resampling,
while empirically competitive, lacks finite-sample guarantees and incurs
substantial computational cost. Our results demonstrate that quantifying
musical style similarity with statistical confidence is both feasible and
principled, opening avenues for reliable downstream tasks such as
AI-generated music filtering and quality assessment.
\end{abstract}

% --- Sections ---
%% ===========================================================================
%%  1. Introduction
%%  Quantifying with Confidence: Conformal Prediction for Reliable Musical
%%  Style Similarity Ranking
%% ===========================================================================

\section{Introduction}

Musical style similarity assessment is a fundamental problem in music
information retrieval (MIR) and a critical component in the evaluation of
AI-generated music. Whether curating playlists, recommending tracks, or
filtering outputs from generative models such as MusicGen~\cite{copet2023musicgen}
and AudioLDM~2~\cite{liu2023audioldm2}, practitioners need to determine how closely
a piece of music aligns with a target stylistic reference. The reliability of
these similarity judgments directly impacts the quality of downstream
applications, yet the methods used to produce them have largely remained
point-estimate heuristics with no principled uncertainty quantification.

Recent advances in self-supervised learning (SSL) have yielded powerful audio
representation models capable of capturing high-level musical attributes
including timbre, harmony, rhythm, and genre. Among these, the Music
undERstanding model with large-scale self-supervised Training
(MERT)~\cite{li2023mert}---a 330M-parameter Transformer pre-trained on 160,000
hours of music data---has demonstrated state-of-the-art performance across a
wide range of MIR tasks. By extracting embeddings from MERT's deep hidden
representations, one can compute pairwise cosine similarity between tracks
and obtain a single scalar score reflecting their stylistic proximity.
However, this similarity score is a point estimate: it conveys no information
about its own reliability, does not account for the inherent ambiguity of
musical style boundaries, and offers no statistical guarantees. In
safety-critical or curation-sensitive settings---such as filtering
AI-generated music for stylistic compliance or constructing evaluation
benchmarks for generative models---the absence of uncertainty quantification
constitutes a significant limitation.

Conformal Prediction (CP)~\cite{vovk2005algorithmic,shafer2008tutorial}
provides a distribution-free, finite-sample framework for constructing
prediction sets with rigorous coverage guarantees. Given a user-specified
significance level $\alpha$, CP produces a prediction set $\Gamma(x)$ that
is guaranteed to contain the correct label with probability at least
$1 - \alpha$, under only the assumption that calibration and test data are
exchangeable. Unlike Bayesian methods, CP makes no parametric assumptions
about the underlying distribution; unlike bootstrap approaches, it offers
provable finite-sample validity rather than asymptotic approximations.
These properties make CP particularly well-suited to subjective perceptual
tasks such as musical style assessment, where the score distribution is
inherently heavy-tailed and resists simple parametric modeling.

In this paper, we propose the first integration of conformal prediction with
self-supervised audio embeddings for reliable musical style similarity
ranking. Our framework extracts 1024-dimensional embeddings from
MERT-v1-330M for each audio track, constructs per-genre prototype centroids
from a reference set, and defines a nonconformity score based on the cosine
distance between a query embedding and each genre centroid. Using a held-out
calibration set, we compute empirical $p$-values and construct prediction
sets with distribution-free coverage guarantees. We evaluate our method on
the GTZAN dataset~\cite{tzanetakis2002gtzan} (999 tracks across 10 genres
after excluding one corrupted file), comparing against bootstrap resampling
and Gaussian approximation baselines under a stratified 5-fold cross-validation
protocol with a 60/20/20 (reference/calibration/test) split.

We make the following contributions:

\begin{enumerate}
    \item We propose the first conformal prediction framework for musical style
    similarity assessment, bridging the gap between self-supervised audio
    representation learning and statistically rigorous uncertainty
    quantification (Section~\ref{sec:method}).

    \item We demonstrate empirically that our CP approach achieves empirical
    coverage tightly matching nominal levels: $0.993 \pm 0.010$ at
    $\alpha = 0.01$ (target $0.99$), $0.955 \pm 0.021$ at $\alpha = 0.05$
    (target $0.95$), and $0.910 \pm 0.028$ at $\alpha = 0.10$
    (target $0.90$). In contrast, the Gaussian approximation baseline
    yields only $0.965$ coverage at $\alpha = 0.01$, confirming that the
    normality assumption fails in the heavy-tailed regime characteristic
    of perceptual similarity tasks (Section~\ref{sec:experiments}).

    \item We provide a systematic comparison of CP against bootstrap
    resampling and Gaussian approximation baselines, highlighting that
    only CP offers a distribution-free, finite-sample coverage guarantee---a
    critical property when the calibration set is small or the underlying
    distribution is non-Gaussian (Section~\ref{sec:experiments}).
\end{enumerate}

The remainder of this paper is organized as follows. Section~\ref{sec:related}
reviews related work on self-supervised music representation, conformal
prediction, and uncertainty quantification in MIR. Section~\ref{sec:method}
presents our CP framework, detailing the MERT embedding pipeline, genre
centroid construction, nonconformity score design, and prediction set
formation. Section~\ref{sec:experiments} describes our experimental setup,
compares CP against baselines, and analyzes coverage and set size results.
Section~\ref{sec:discussion} discusses the implications of our findings,
including the theoretical advantages of CP in perceptual tasks and
limitations of the current approach. Section~\ref{sec:conclusion} concludes
with a summary of contributions and directions for future work.

%% ===========================================================================
%%  2. Related Work
%% ===========================================================================

\section{Related Work}
\label{sec:related}

Our work sits at the intersection of three research areas: (i)~trust and
explainability in human-AI interaction, (ii)~self-supervised audio
representation learning for music, and (iii)~conformal prediction for
distribution-free uncertainty quantification. We review each in turn and
identify the gap that our framework fills.

\subsection{Trust and Explainability in Human-AI Interaction}

A central challenge in human-AI interaction is that users cannot assess the
reliability of AI outputs when only point predictions are
provided~\cite{amershi2019guidelines}. This trust gap is particularly acute in
creative and perceptual domains, where ground truth is ambiguous and system
errors may go unnoticed by non-experts. Miller~\cite{miller2019explanation}
argues that explanations in AI should be contrastive, selective, and social---in
other words, they should help users understand \emph{why} a particular output
was produced rather than \emph{how} the model works internally.

Existing approaches to communicating model confidence fall into two broad
categories. \emph{Post-hoc explanation} methods, such as
LIME~\cite{ribeiro2016lime}, highlight input features that influenced a
prediction but provide no statistical guarantee about the prediction's
correctness. \emph{Bayesian uncertainty} techniques, including Monte Carlo
dropout~\cite{gal2016dropout} and deep ensembles, produce variance estimates
that are difficult for non-technical users to interpret and offer only
asymptotic calibration.

Our work proposes a third path: instead of explaining \emph{why} a model chose a
label, we communicate \emph{how many} labels could plausibly be correct. The
prediction set size---a single integer---provides an immediately interpretable
confidence signal that requires no understanding of probability, entropy, or
model architecture. This aligns with the principle articulated by Amershi et
al.~\cite{amershi2019guidelines} that AI systems should ``convey the
consequences of user actions'' and ``support efficient dismissal'' of
suggestions.

\subsection{Music AI and Creative Workflows}

The adoption of AI in music tools has grown rapidly, from style
transfer~\cite{huang2020creative} to generative composition~\cite{copet2023musicgen}
to automated playlist curation. In these applications, genre classification
serves as a fundamental building block for downstream tasks such as filtering,
tagging, and recommendation. However, music genre is inherently subjective and
multidimensional---a single track may legitimately belong to multiple genres,
and annotators frequently disagree~\cite{tzanetakis2002gtzan}. Traditional
single-label classifiers force this ambiguity into a single output, misleading
users about the certainty of the classification.

For generative music specifically, evaluation remains an open
challenge~\cite{yang2018evaluation}. Metrics such as the Fr\'{e}chet Audio
Distance (FAD)~\cite{kilgour2019fad} measure distribution-level similarity but
provide no per-sample feedback. A producer cannot ask ``is \emph{this
particular} generated track convincing as jazz?'' and receive an answer with
quantified confidence. Our prediction sets address this need by providing
instance-level reliability information.

\subsection{Self-Supervised Audio Representations}

Self-supervised learning (SSL) has become the dominant paradigm for extracting
features from raw audio. In the speech domain, Wav2Vec~2.0~\cite{baevski2020wav2vec2}
and HuBERT~\cite{hsu2021hubert} demonstrated that learned representations
surpass hand-crafted features on a wide range of tasks. For music,
MERT~\cite{li2023mert} extends this approach: a 24-layer Transformer
pre-trained on 160,000 hours of music via masked language modeling with dual
teacher supervision (acoustic RVQ-VAE and musical CQT). MERT produces
1024-dimensional embeddings that encode both timbral and harmonic information,
and has been shown to support strong performance across MIR benchmarks when
paired with task-specific models.

In this work, we use MERT-v1-330M purely as a frozen feature extractor,
extracting embeddings from the last hidden layer via mean-pooling over time. We
then train a lightweight Linear Probe on these embeddings, keeping the
underlying representation fixed. This design choice keeps training fast and
reproducible while leveraging the rich acoustic knowledge embedded in MERT's
pre-trained weights. For a detailed comparison of audio SSL models and
alternative multimodal approaches such as CLAP, we refer the reader to
comprehensive surveys in the MIR literature.

\subsection{Conformal Prediction as a User-Facing Interface}

Conformal prediction (CP)~\cite{vovk2005algorithmic,shafer2008tutorial} is a
distribution-free, finite-sample framework for constructing prediction sets with
formal coverage guarantees: given a significance level $\alpha$, CP produces a
set $\Gamma(x)$ that contains the true label with probability at least
$1-\alpha$, under only the assumption that calibration and test data are
exchangeable. A modern introduction is provided by Angelopoulos and
Bates~\cite{angelopoulos2023gentle}.

CP has primarily been studied as a statistical methodology for machine learning
practitioners. Applications include regression~\cite{romano2019conformalized},
image classification with class-conditional
sets~\cite{angelopoulos2021uncertainty}, and software libraries for practical
deployment~\cite{cordier2023mapie}. However, to our knowledge, no prior work has
explicitly positioned CP as a \emph{user interface layer}---a mechanism through
which AI systems communicate their certainty to non-expert users in an
interpretable, actionable format. This is the perspective we adopt.

From an HCI standpoint, CP has several attractive properties as an interface
primitive:
\begin{enumerate}[leftmargin=*,itemsep=0pt]
    \item \textbf{Set size as a natural confidence indicator.} Unlike
    probabilities or entropy values, the cardinality of a prediction set is
    trivially interpretable: smaller means more confident, larger means more
    uncertain. This maps to the long-established finding that humans reason more
    reliably about categorical sets than about continuous probability
    values~\cite{miller2019explanation}.
    \item \textbf{$\alpha$ as a direct user control.} The significance level
    functions as a slider that users can adjust according to their risk
    tolerance, without needing to understand the underlying calibration
    mechanism.
    \item \textbf{Distribution-free guarantees.} The coverage guarantee holds
    regardless of model architecture, embedding space geometry, or score
    distribution shape, making the interface behavior predictable and uniform
    across different underlying systems.
\end{enumerate}

Our contribution is to realize this interface perspective concretely for music
genre classification, demonstrating that CP not only provides theoretical
guarantees but also yields practically useful prediction sets that can inform
user decision-making in creative workflows.

%% ===========================================================================
%%  3. Methodology
%% ===========================================================================

\section{Methodology}
\label{sec:methodology}

Given a dataset of $N$ audio tracks with genre labels $y_i \in \mathcal{Y} =
\{1, \dots, K\}$, our goal is to produce, for any test track $x$, a
\emph{prediction set} $\Gamma_{\alpha}(x) \subseteq \mathcal{Y}$ satisfying
$\mathbb{P}(y_{\text{test}} \in \Gamma_{\alpha}(x_{\text{test}})) \geq
1 - \alpha$, where $\alpha$ is chosen by the user. We achieve this through three
stages.

\subsection{Embedding Extraction}

We use MERT-v1-330M~\cite{li2023mert}, a 24-layer Transformer pre-trained on
160,000 hours of music, as a frozen feature extractor. Audio is resampled to
24\,kHz, converted to mono, split into 10\,s windows with 50\% overlap, and
passed through MERT. Hidden states from the last layer are mean-pooled over
time, averaged across windows, and $\ell_2$-normalized:
\begin{equation}
  \mathbf{e}_x = \frac{1}{S} \sum_{s=1}^{S}
  \frac{1}{T_s} \sum_{t=1}^{T_s} \mathbf{h}_L^{(s)}[t] \in \mathbb{R}^{1024}.
  \label{eq:embedding}
\end{equation}

\subsection{Learning Discriminative Prototypes}

We train a Linear Probe---a single linear layer $\mathbf{W} \in \mathbb{R}^{K
\times d}$ with bias $\mathbf{b}$---on a reference set
$\mathcal{D}_{\text{ref}}$ (60\% of data) by minimizing cross-entropy with L2
regularization (Adam, $\text{lr}=10^{-3}$, $\text{weight\_decay}=10^{-2}$, 100
epochs):
\begin{equation}
  \mathcal{L}(\mathbf{W}, \mathbf{b}) = -\sum_{i \in \mathcal{D}_{\text{ref}}}
  \log \frac{\exp(\mathbf{W}_{y_i} \cdot \mathbf{e}_i + b_{y_i})}
  {\sum_{k=1}^{K} \exp(\mathbf{W}_k \cdot \mathbf{e}_i + b_k)}.
  \label{eq:lp_loss}
\end{equation}
The trained row vectors $\mathbf{W}_k$ serve as discriminative genre prototypes
that actively separate class boundaries, in contrast to simple centroid
averaging which passively inherits class overlap.

\subsection{Conformal Calibration of Prediction Sets}

We adopt inductive (split) CP with a 60/20/20 stratified split into reference,
calibration ($\mathcal{D}_{\text{cal}}$), and test sets. The nonconformity score
for an embedding $\mathbf{e}_x$ under candidate genre $k$ is the cosine distance
to prototype $\mathbf{W}_k$:
\begin{equation}
  A(\mathbf{e}_x, k) = 1 - \cos(\mathbf{e}_x, \mathbf{W}_k).
  \label{eq:lp_score}
\end{equation}
We use cosine distance (rather than softmax residual) for its geometric
interpretability: each $\mathbf{W}_k$ is a direction in embedding space, and
$A$ measures angular deviation independent of other classes.

Let $\mathcal{A}_{\text{cal}} = \{ A(\mathbf{e}_{x_i}, y_i) \mid (x_i, y_i) \in
\mathcal{D}_{\text{cal}} \}$ be the set of calibration nonconformity scores
computed against true labels. The empirical $p$-value for a test embedding under
candidate genre $k$ is the fraction of calibration scores that are at least as
large:
\begin{equation}
  p(\mathbf{e}_x, k) = \frac{
    |\{ A_i \in \mathcal{A}_{\text{cal}} : A_i \geq A(\mathbf{e}_x, k) \}| + 1
  }{
    |\mathcal{D}_{\text{cal}}| + 1
  }.
  \label{eq:pvalue}
\end{equation}
The prediction set includes all genres whose $p$-value exceeds $\alpha$:
\begin{equation}
  \Gamma_{\alpha}(x) = \{ k \in \mathcal{Y} : p(\mathbf{e}_x, k) > \alpha \}.
  \label{eq:prediction_set}
\end{equation}
This construction guarantees $\mathbb{P}(y_{\text{test}} \in
\Gamma_{\alpha}(x_{\text{test}})) \geq 1 - \alpha$ under exchangeability. From
the user's perspective, $|\Gamma_{\alpha}(x)|$ serves as an intuitive confidence
indicator (smaller = more confident), and $\alpha$ serves as a trust slider
adjustable by task.

We compare three calibration strategies sharing the identical Linear Probe
model: (i)~LP-Cos-CP (ours) uses empirical $p$-values with cosine scores;
(ii)~LP-Softmax-Bootstrap uses softmax residual scores with $B=1{,}000$
bootstrap calibration; (iii)~LP-Softmax-Gaussian fits
$\mathcal{N}(\hat{\mu}, \hat{\sigma}^2)$ to calibration scores. Cosine and
softmax scores produce near-identical $p$-values (difference $< 0.005$), so
this comparison isolates calibration strategy. We evaluate empirical coverage
and average set size over 5 stratified folds at $\alpha \in \{0.01, 0.05,
0.10\}$.

%% ===========================================================================
%%  4. Experiments
%%  Quantifying with Confidence: Conformal Prediction for Reliable Musical
%%  Style Similarity Ranking
%% ===========================================================================

\section{Experiments}
\label{sec:experiments}

We evaluate the proposed framework along three dimensions:
(1)~the quality of MERT embeddings for capturing stylistic similarity,
(2)~the coverage and set size characteristics of conformal prediction sets,
and (3)~the comparison of CP against bootstrap and Gaussian approximation
baselines. All experiments use the GTZAN dataset under a stratified 5-fold
cross-validation protocol.

\subsection{Dataset and Preprocessing}

\textbf{GTZAN}~\cite{tzanetakis2002gtzan} is the de facto standard benchmark
for music genre classification and a common testbed for music information
retrieval research. It consists of 1{,}000~audio tracks evenly distributed
across 10~genres: blues, classical, country, disco, hip-hop, jazz, metal,
pop, reggae, and rock. Each track is a 30-second monophonic WAV file sampled
at 22.05\,kHz.

During data validation, we identified one corrupted file
(\texttt{jazz.00054.wav}) that could not be successfully decoded, leaving
$N = 999$ tracks for the experiments. The remaining jazz class has 99 tracks;
all other classes have 100~tracks. This minor imbalance is accounted for by
our stratified splitting procedure and is negligible for the coverage analysis.

\textbf{Preprocessing.} All audio tracks are resampled to MERT's native
sampling rate of 24\,kHz, converted to mono via channel averaging, and
processed through the embedding extraction pipeline described in
Section~\ref{sec:method}. The output is a matrix
$\mathbf{E} \in \mathbb{R}^{999 \times 1024}$ of $\ell_2$-normalized
embeddings. Extraction of all embeddings requires approximately 9~minutes on
an NVIDIA RTX~5070~GPU (8\,GB~VRAM).

\subsection{Experimental Setup}

We adopt a stratified 5-fold protocol that preserves per-genre proportions
across all subsets. For each fold, the data is partitioned into three disjoint
splits:
\begin{itemize}[leftmargin=*,itemsep=1pt]
    \item \textbf{Reference set} ($60\%$): used to compute per-genre prototype
    centroids $\boldsymbol{\mu}_k$ via Equation~\ref{eq:centroid}.
    \item \textbf{Calibration set} ($20\%$): used to compute the empirical
    distribution of nonconformity scores $\mathcal{A}_{\text{cal}}$ and
    calibrate prediction thresholds.
    \item \textbf{Test set} ($20\%$): held out for evaluation; never used
    during centroid construction or calibration.
\end{itemize}
For each test point $x$, we compute nonconformity scores
$A(\mathbf{e}_x, k)$ with respect to all $K = 10$ genre centroids, derive
$p$-values via Equation~\ref{eq:pvalue}, and construct prediction sets
$\Gamma_{\alpha}(x)$ via Equation~\ref{eq:prediction_set}. Results are
aggregated across all five folds to produce means and standard deviations.

We evaluate at three significance levels $\alpha \in \{0.01, 0.05, 0.10\}$,
corresponding to confidence targets of $99\%$, $95\%$, and $90\%$,
respectively. Random seeds are fixed (base seed 42, incremented by fold
index) to ensure reproducibility.

\subsection{Embedding Quality Verification}

Before evaluating the conformal prediction framework, we assess whether
MERT embeddings capture genre-discriminative structure in the GTZAN dataset.
We compute two quantities: (i)~\emph{intra-class cosine similarity}---the
mean pairwise cosine similarity between embeddings of tracks belonging to
the same genre; and (ii)~\emph{inter-class cosine similarity}---the mean
cosine similarity between embeddings of tracks belonging to different genres.

The intra-class similarity averaged across all 10 genres is $0.917$, while
the inter-class similarity is $0.873$, yielding a separation margin of
$\Delta = 0.044$. The positive margin confirms that tracks from the same
genre are systematically closer in the MERT embedding space than tracks from
different genres. A two-dimensional PCA projection of the embedding matrix
(see Figure~\ref{fig:pca} in Appendix~\ref{sec:appendix}) further reveals
clear clustering tendencies, with genres such as \emph{classical} and
\emph{metal} forming tight, well-separated groups and more ambiguous genres
such as \emph{rock} and \emph{country} occupying overlapping regions---a
pattern consistent with known perceptual ambiguities in musical genre
taxonomy.

We detect no NaN or infinite values in the embedding matrix, confirming that
the MERT extraction pipeline is numerically stable.

\subsection{Coverage Analysis}

Table~\ref{tab:coverage_results} reports the empirical coverage of CP,
Bootstrap, and Gaussian methods averaged over the five cross-validation folds.

\begin{table}[htbp]
\centering
\caption{Empirical coverage and average prediction set size across three
significance levels (\textbf{bold} = best coverage). Coverage $\geq$ nominal
is highlighted as matching or exceeding the $1-\alpha$ target. Results
aggregated over 5 stratified folds (mean $\pm$ std).
}
\label{tab:coverage_results}
\begin{tabular}{c c c c c}
\toprule
$\alpha$ & Nominal $1-\alpha$ & CP & Bootstrap & Gaussian \\
\midrule
\addlinespace[2pt]
\multirow{2}{*}{0.01} & \multirow{2}{*}{0.99}
    & \textbf{0.993 $\pm$ 0.010} & 0.987 $\pm$ 0.012 & 0.965 $\pm$ 0.020 \\
    & & \hspace{2pt}(set size: 8.51) & (8.11) & (7.05) \\[4pt]
\multirow{2}{*}{0.05} & \multirow{2}{*}{0.95}
    & \textbf{0.955 $\pm$ 0.021} & 0.950 $\pm$ 0.023 & 0.932 $\pm$ 0.033 \\
    & & \hspace{2pt}(set size: 6.69) & (6.59) & (6.03) \\[4pt]
\multirow{2}{*}{0.10} & \multirow{2}{*}{0.90}
    & \textbf{0.910 $\pm$ 0.028} & \textbf{0.910 $\pm$ 0.029} & 0.907 $\pm$ 0.031 \\
    & & \hspace{2pt}(set size: 5.40) & (5.39) & (5.35) \\
\bottomrule
\end{tabular}
\end{table}

\textbf{Conformal prediction meets its coverage guarantee.} At
$\alpha = 0.01$, CP achieves coverage of $0.993$, exceeding the nominal
$0.99$ target. At $\alpha = 0.05$, coverage is $0.955$, and at $\alpha =
0.10$, it is $0.910$---both tightly matching their respective nominal levels.
The standard deviations (0.010--0.028) reflect the fold-to-fold variability
expected from a 999-track dataset, and are consistent with the
$\mathcal{O}(1/\sqrt{n})$ scaling of empirical coverage estimates.

\textbf{Gaussian approximation fails at high confidence.} The Gaussian
baseline achieves only $0.965$ coverage at $\alpha = 0.01$---a shortfall of
$2.5$~percentage~points below the $0.99$ nominal target. This confirms that
the nonconformity score distribution in the MERT embedding space deviates
from normality, exhibiting heavier tails than a Gaussian distribution. At
more lenient significance levels ($\alpha = 0.05, 0.10$), the Gaussian
approximation approaches acceptable coverage, consistent with the asymptotic
robustness of Gaussian quantiles in the central regime. However, the
deficiency at $\alpha = 0.01$ is precisely the regime most relevant for
high-stakes filtering and safety-critical applications.

\textbf{Bootstrap performs comparably to CP but without guarantees.} The
bootstrap baseline achieves coverage of $0.987$ at $\alpha = 0.01$, slightly
below the nominal target, and $0.950$ and $0.910$ at $\alpha = 0.05$ and
$\alpha = 0.10$, respectively. While its aggregate performance closely
tracks CP on the GTZAN dataset, bootstrap offers no finite-sample coverage
guarantee and requires $B = 1{,}000$ resamples of the calibration set,
incurring roughly $1{,}000\times$ the computational cost of CP calibration.

\subsection{Prediction Set Size Analysis}

The second row for each $\alpha$ in Table~\ref{tab:coverage_results} reports
the average prediction set size. The key observations are:

\begin{itemize}[leftmargin=*,itemsep=1pt]
    \item \textbf{Set size decreases with $\alpha$, as expected.} As the
    significance level relaxes ($\alpha: 0.01 \to 0.10$), the
    $(1 - \alpha)$-quantile of calibration scores decreases, admitting fewer
    genres into the prediction set. CP set sizes shrink from $8.51$ (at
    $\alpha = 0.01$) to $5.40$ (at $\alpha = 0.10$), reflecting increasing
    permissiveness of the inclusion criterion.

    \item \textbf{Gaussian achieves smaller sets through undercoverage.} The
    Gaussian baseline consistently produces the smallest prediction sets
    (e.g., $7.05$ vs.\ CP's $8.51$ at $\alpha = 0.01$), but this apparent
    advantage is spurious: it stems from an overly optimistic threshold
    derived from the normality assumption, which excludes true labels that
    fall in the heavy tail. The smaller set size is not a sign of higher
    precision but of \emph{overconfidence}---a well-known pitfall of
    parametric methods on non-Gaussian data.

    \item \textbf{CP achieves the best coverage--size trade-off.} At
    $\alpha = 0.01$, CP's set size of $8.51$ is approximately $1.46$ larger
    than the Gaussian baseline's $7.05$, yet CP provides valid coverage
    ($0.993 \geq 0.99$) while Gaussian does not ($0.965 < 0.99$). This
    demonstrates that CP pays an appropriate ``uncertainty premium'' in set
    size to deliver its distribution-free coverage guarantee---a premium that
    is both principled and empirically necessary.
\end{itemize}

The distribution of set sizes across individual test tracks (visualized in
Figure~\ref{fig:set_size_boxplot} in Appendix~\ref{sec:appendix}) reveals
substantial heterogeneity: while most tracks receive prediction sets
containing 4--9 genres, a small number of tracks are assigned either very
small sets (1--2 genres, indicating high confidence) or the full set of 10
genres (indicating maximal ambiguity). Empty prediction sets are rare
($<1\%$ across all folds and $\alpha$ levels), consistent with the fact that
GTZAN tracks are drawn from the same distribution as the calibration data.

\subsection{Cross-Dataset Validation on FMA}
\label{sec:fma_experiment}

To assess whether the conformal prediction framework generalizes beyond
GTZAN, we conduct a parallel experiment on the Free Music Archive (FMA) Small
dataset~\cite{defferrard2017fma}---a substantially larger and more diverse
collection of full-length Creative Commons-licensed music. FMA Small contains
8{,}000 tracks (7{,}994 after excluding corrupted files) spanning eight
top-level genres: Electronic, Experimental, Folk, Hip-Hop, Instrumental,
International, Pop, and Rock. Tracks are variable-length MP3 files (most 30--180\,s)
encoding a broader range of musical styles and recording conditions than
GTZAN's curated 30-second excerpts.

\textbf{Embedding quality.} Following the same MERT extraction pipeline
(Section~\ref{sec:method}), we obtain a $7{,}994 \times 1024$ embedding matrix
for FMA Small. The intra-class cosine similarity is $0.847$, the inter-class
similarity is $0.821$, yielding a separation margin of $\Delta = 0.026$. This
is substantially narrower than GTZAN's $\Delta = 0.044$---approximately
$41\%$ weaker separation---consistent with the larger, noisier, and
stylistically more heterogeneous composition of the FMA dataset. The positive
margin nonetheless confirms that MERT embeddings exhibit meaningful
genre-discriminative structure.

\textbf{CP results on FMA.} Table~\ref{tab:fma_results} reports the empirical
coverage and set size for all three methods on FMA, aggregated over five
stratified folds.

\begin{table}[htbp]
\centering
\caption{Empirical coverage and average prediction set size on the FMA Small
dataset (\textbf{bold} = best coverage). Results aggregated over 5 stratified
folds (mean $\pm$ std). FMA contains more tracks ($\sim$8{,}000) and exhibits
weaker class separation ($\Delta = 0.026$) than GTZAN.
}
\label{tab:fma_results}
\begin{tabular}{c c c c c}
\toprule
$\alpha$ & Nominal $1-\alpha$ & CP & Bootstrap & Gaussian \\
\midrule
\addlinespace[2pt]
\multirow{2}{*}{0.01} & \multirow{2}{*}{0.99}
    & \textbf{0.988 $\pm$ 0.004} & 0.987 $\pm$ 0.004 & 0.966 $\pm$ 0.003 \\
    & & \hspace{2pt}(set size: 7.75) & (7.72) & (7.22) \\[4pt]
\multirow{2}{*}{0.05} & \multirow{2}{*}{0.95}
    & \textbf{0.948 $\pm$ 0.004} & 0.947 $\pm$ 0.004 & 0.934 $\pm$ 0.005 \\
    & & \hspace{2pt}(set size: 6.85) & (6.83) & (6.59) \\[4pt]
\multirow{2}{*}{0.10} & \multirow{2}{*}{0.90}
    & 0.898 $\pm$ 0.006 & 0.896 $\pm$ 0.006 & \textbf{0.903 $\pm$ 0.005} \\
    & & \hspace{2pt}(set size: 6.03) & (6.02) & (6.10) \\
\bottomrule
\end{tabular}
\end{table}

Several observations are noteworthy. First, CP continues to provide coverage
closely matching nominal targets despite the weaker embedding space---$0.988$
at $\alpha = 0.01$ (target $0.99$) and $0.948$ at $\alpha = 0.05$ (target
$0.95$). This demonstrates that the distribution-free guarantee is robust to
the quality of the underlying representation: CP adapts its threshold to the
empirical score distribution regardless of the discriminability of the
embedding space.

Second, the standard deviations on FMA ($\sim$0.003--0.006) are substantially
smaller than on GTZAN ($\sim$0.010--0.033), reflecting the nearly $8\times$
larger sample size and the corresponding reduction in fold-to-fold variance.

Third, the Gaussian baseline at $\alpha = 0.10$ achieves slightly higher
coverage ($0.903$) than CP ($0.898$), an inversion of the pattern observed on
GTZAN. This occurs because the Gaussian quantile is more conservative than CP's
empirical quantile at this particular significance level for the FMA score
distribution. However, at $\alpha = 0.01$---the high-confidence regime most
critical for downstream applications---Gaussian undercovers by $2.4$ percentage
points ($0.966$ vs.\ target $0.99$), consistently failing where guarantees
matter most.

Finally, the prediction set sizes on FMA are generally smaller than on GTZAN
(e.g., $7.75$ vs.\ $8.51$ at $\alpha = 0.01$), despite the weaker class
separation. This counterintuitive result can be explained by two factors:
(1)~the substantially larger calibration set ($\sim$1{,}600 tracks vs.\
$\sim$200 on GTZAN) provides a more precise empirical quantile estimate,
enabling tighter thresholds; and (2)~FMA's eight genres present a simpler
prediction task ($K = 8$) than GTZAN's ten ($K = 10$), reducing the
theoretical upper bound on set size.

\subsection{Summary of Experimental Findings}

Our experiments yield four primary takeaways:

\begin{enumerate}[leftmargin=*,itemsep=1pt]
    \item \textbf{MERT embeddings are genre-discriminative.} On GTZAN, the
    intra-class similarity ($0.917$) meaningfully exceeds the inter-class
    similarity ($0.873$), confirming that MERT captures stylistic structure.
    On FMA, the weaker separation ($\Delta = 0.026$) reflects the dataset's
    greater diversity but remains above chance, with CP adapting to the
    reduced discriminability without compromising its coverage guarantee.

    \item \textbf{Conformal prediction provides valid, tight coverage across
    datasets.} At all three significance levels on both GTZAN and FMA, CP
    achieves empirical coverage at or near the nominal target, with only
    moderate set sizes. On FMA, CP achieves $0.988$ coverage at $\alpha=0.01$
    despite the substantially weaker embedding space ($\Delta=0.026$ vs.\
    GTZAN's $\Delta=0.044$), demonstrating that the distribution-free
    guarantee is representation-agnostic.

    \item \textbf{Parametric baselines fail where CP succeeds.} The Gaussian
    approximation baseline undercovers at $\alpha = 0.01$ on \emph{both}
    datasets---$0.965$ on GTZAN and $0.966$ on FMA (vs.\ target $0.99$).
    The pervasiveness of this failure across datasets with markedly different
    score distributions underscores that the normality assumption is
    fundamentally inappropriate for perceptual similarity tasks. Bootstrap
    resampling, while empirically competitive, lacks the finite-sample
    validity guarantee that distinguishes CP.

    \item \textbf{Larger calibration sets yield tighter prediction sets.} The
    FMA calibration set ($\sim$1{,}600 tracks) is approximately $8\times$
    larger than GTZAN's ($\sim$200 tracks), resulting in prediction sets that
    are consistently smaller on FMA ($7.75$ vs.\ $8.51$ at $\alpha=0.01$)
    despite weaker class separation. This illustrates a key practical
    principle: deploying CP with a sufficiently large calibration corpus
    narrows the gap between the coverage guarantee and practical
    informativeness.
\end{enumerate}

%% ===========================================================================
%%  5. Discussion
%% ===========================================================================

\section{Discussion}
\label{sec:discussion}

Our results demonstrate that CP can function as a user-facing interface layer
between AI music models and their human users. We distill three design
principles and discuss limitations.

\subsection{Design Implications}

\textbf{Prediction set size is an intuitive confidence indicator.} 
Existing approaches to communicating model uncertainty often rely on calibrated confidence displays (e.g., exposing raw probability distributions) or selective prediction (e.g., abstaining entirely when uncertain)~\cite{bhatt2021uncertainty,geifman2017selective}. However, continuous probability distributions are often difficult for lay users to interpret correctly~\cite{padilla2021uncertainty}, and pure selective prediction can disrupt workflows by offering no guidance. CP prediction sets provide a powerful middle ground: set cardinality translates complex model uncertainty into a discrete, actionable format without requiring statistical training. $|\Gamma| = 1$ signals certainty, while $|\Gamma| \ge 4$ signals ambiguity. This aligns with Miller's observation that humans reason more reliably about categorical sets than continuous probabilities~\cite{miller2019explanation}. A UI showing a track's valid prediction set enables users to calibrate trust at a glance.

\textbf{$\alpha$ provides a direct user control.} The significance level maps
to an answerable question: ``What risk of missing the true genre am I willing to
accept?'' This enables a spectrum from conservative mode ($\alpha = 0.01$, 6--7
candidates) for high-stakes curation to exploratory mode ($\alpha = 0.10$, 2--3
candidates) for casual browsing, all within the same system.

\textbf{Distribution-free calibration is essential for predictable interface
behavior.} The Gaussian baseline's mass overcoverage and Bootstrap's marginal
undercoverage under the identical prediction model demonstrate that parametric
calibration assumptions produce unreliable user-facing behavior. Only empirical
$p$-value calibration consistently respects the coverage promise across datasets
and $\alpha$ levels.

\subsection{Limitations and Future Work}

\textbf{Exchangeability requirement.} The coverage guarantee assumes calibration
and test data are exchangeable. When deploying on AI-generated music or novel
distributions, this may be violated. Integrating novelty detection that alerts
users when a track falls outside the calibration distribution is a natural
extension.

\textbf{Static prototypes.} The Linear Probe produces fixed prototypes that
cannot adapt to personal genre ontologies or evolving styles. Personalizing
prototypes via user feedback (genre corrections, listening history) could
improve relevance while preserving the CP guarantee.

\textbf{User studies needed.} The present work establishes algorithmic
foundations. A critical next step is evaluating how prediction sets affect user
trust calibration, decision speed, and task accuracy in realistic music curation
scenarios. We believe the quantitative results---valid coverage, compact sets,
and controllable behavior---provide compelling motivation for such studies.

Beyond music, prediction sets are a general-purpose interface primitive. We
encourage investigation of CP as a user-facing layer in other perceptual domains
where single-label predictions fail to communicate the ambiguity inherent in
human perception.

%% ===========================================================================
%%  6. Conclusion
%% ===========================================================================

\section{Conclusion}
\label{sec:conclusion}

We presented LP-Cos-CP, a framework that replaces opaque single-genre
predictions with interpretable, user-controllable prediction sets. On GTZAN and
FMA Small, our method achieves empirical coverage tightly matching nominal
targets while reducing set sizes by up to 55\% compared to centroid
baselines. The $\alpha$ parameter functions as a principled trust slider,
enabling users to configure the reliability--informativeness trade-off for their
specific task.

Our comparison of calibration strategies reveals that distribution-free
calibration is not merely a statistical nicety but a practical requirement for
predictable user-facing behavior: parametric assumptions (Gaussian,
Bootstrap) produce interface behavior that is either uninformative or
untrustworthy. CP's empirical $p$-value calibration remains robust across
datasets and score distributions.

The prediction set provides a template for bridging the gap between ML
performance and human usability. Future work includes user studies evaluating
prediction sets in interactive music tools, novelty detection for
out-of-distribution tracks, and prototype personalization through user feedback.
We believe this interface paradigm extends naturally to other perceptual
classification domains.


% --- Figures ---
\clearpage
\section*{Figures}
\label{sec:appendix}

\begin{figure}[H]
    \centering
    \includegraphics[width=0.85\textwidth]{../outputs/figures/coverage_vs_alpha.png}
    \caption{Coverage vs.\ significance level $\alpha$. The dashed diagonal
    represents the nominal $1-\alpha$ target. Error bars show $\pm 1$ standard
    deviation across 5 stratified folds.}
    \label{fig:coverage_vs_alpha}
\end{figure}

\begin{figure}[H]
    \centering
    \includegraphics[width=0.85\textwidth]{../outputs/figures/coverage_deviation.png}
    \caption{Coverage deviation (empirical $-$ nominal) across three significance
    levels. Bars above zero indicate overcoverage; bars below zero indicate
    undercoverage.}
    \label{fig:coverage_deviation}
\end{figure}

\begin{figure}[H]
    \centering
    \includegraphics[width=0.95\textwidth]{../outputs/figures/set_size_boxplot.png}
    \caption{Prediction set size distribution across significance levels.
    Boxes show quartiles; whiskers show range excluding outliers; diamonds mark
    means.}
    \label{fig:set_size_boxplot}
\end{figure}

\begin{figure}[H]
    \centering
    \includegraphics[width=0.85\textwidth]{../outputs/figures/pca_embeddings.png}
    \caption{Two-dimensional PCA projection of the 1024-dimensional MERT
    embeddings. Each point corresponds to one GTZAN track; colors indicate
    genre labels.}
    \label{fig:pca}
\end{figure}

\begin{figure}[H]
    \centering
    \includegraphics[width=0.85\textwidth]{../outputs/figures/fma_coverage_comparison.png}
    \caption{CP coverage on GTZAN vs.\ FMA Small. Despite the substantially
    weaker class separation on FMA ($\Delta = 0.026$ vs.\ GTZAN's $0.044$),
    CP maintains near-nominal coverage on both datasets. Error bars show
    $\pm 1$ standard deviation across 5 stratified folds.}
    \label{fig:fma_coverage}
\end{figure}

\begin{figure}[H]
    \centering
    \includegraphics[width=0.85\textwidth]{../outputs/figures/fma_set_size_comparison.png}
    \caption{CP average prediction set size on GTZAN vs.\ FMA Small. Despite
    fewer genres ($K=8$ vs.\ $K=10$) and weaker class separation, FMA produces
    smaller prediction sets at $\alpha = 0.01$ due to its substantially larger
    calibration corpus ($\sim$1{,}600 vs.\ $\sim$200 tracks).}
    \label{fig:fma_set_size}
\end{figure}

\begin{figure}[H]
    \centering
    \includegraphics[width=\textwidth]{../outputs/figures/fma_full_results.png}
    \caption{Full results on FMA Small. (Left) Coverage vs.\ $\alpha$ for all
    three methods. (Right) Coverage deviation from nominal. All methods
    perform comparably at $\alpha = 0.10$, but Gaussian undercovers at
    $\alpha = 0.01$, consistent with the GTZAN findings.}
    \label{fig:fma_full}
\end{figure}

\begin{figure}[H]
    \centering
    \fbox{\parbox{0.8\textwidth}{\centering\vspace{3cm}
    \textbf{[Schematic: Pipeline Overview]}\\[4pt]
    Audio Input $\to$ MERT Encoder $\to$ Embedding $\mathbf{e}_x$
    $\to$ Genre Centroids $\to$ Nonconformity Score
    $\to$ Calibration $\to$ Prediction Set $\Gamma_{\alpha}(x)$
    \vspace{3cm}}}
    \caption{Schematic overview of the proposed framework. Raw audio is
    processed through MERT to extract fixed-dimensional embeddings. Genre
    prototype centroids are computed from a reference set. Nonconformity
    scores for each query track are calibrated against a held-out calibration
    set to produce prediction sets with coverage guarantees.}
    \label{fig:pipeline}
\end{figure}

\clearpage

% --- References ---
\bibliographystyle{ieeetr}
\bibliography{../references/refs}

\end{document}
